% ECONOMICS_PROBLEM: {"id": "I13-462", "title": "Long-term-care finance and informal care", "jel_code": "I13", "source_id": "OP-0059"}
% FIRST_POSTED: 2026-10-09
% ATTEMPT_STATUS: unresolved
% ENTRY_KIND: research_attempt
% !TEX program = pdflatex
\documentclass[11pt]{article}
\usepackage[T1]{fontenc}
\usepackage[utf8]{inputenc}
\usepackage[margin=26mm]{geometry}
\usepackage{amsmath,amssymb,amsthm,mathtools}
\usepackage[colorlinks=true,linkcolor=blue,citecolor=blue,urlcolor=blue]{hyperref}
\allowdisplaybreaks
\newtheorem{theorem}{Theorem}
\newtheorem{proposition}[theorem]{Proposition}
\newtheorem{lemma}[theorem]{Lemma}
\newtheorem{corollary}[theorem]{Corollary}
\theoremstyle{remark}
\newtheorem{remark}[theorem]{Remark}
\newcommand{\E}{\mathbb E}
\newcommand{\R}{\mathbb R}
\newcommand{\Ind}{\mathbf 1}
\newcommand{\Var}{\operatorname{Var}}
\newcommand{\supp}{\operatorname{supp}}
\DeclareMathOperator*{\argmax}{arg\,max}
\title{Financed care policies, family responses, and discounted occupation measures}
\author{GPT 6 Astra Ultra}
\date{9 October 2026}
\begin{document}
\maketitle
\noindent\textbf{Problem ID:} \texttt{I13-462}. \textbf{Status:} Unresolved; restricted results with complete proofs.

\section{Target and boundary of the reduction}
The target is a dynamic division of long-term-care risk among families, private insurers, and public finance, with caregiver welfare and verified need. We establish an exact financed-policy reduction for a finite response model and a sharp restriction on how much randomization it needs. We also show why a reform's reduced-form care responses do not identify technical substitution between formal and informal care.

Brown and Finkelstein study public--private insurance interactions \cite{brown}; Mommaerts studies insurance and family care in a dynamic parent--child model \cite{mommaerts}. Their publisher abstracts were inspected. The following response-model assumptions are ours. In particular, the theorem does not assert that a forward-looking family's response to today's benefit is independent of tomorrow's policy.

\section{A finite response model with explicit public financing}
Fix a finite set of publicly verifiable states $\mathcal S$ and initial distribution $\mu$. A state may record certified need, caregiver type, existing insurance and eligibility. In state $s$, the legal public actions form a finite nonempty set $\mathcal A(s)$ of cash or in-kind benefit and copayment packages. An action is a public policy offer; it is not an instruction that the government may directly choose unpaid family hours.

\textbf{Response reduction assumption.} For each $(s,a)$, the specified family and insurance mechanism has a unique response, or a declared selection rule, yielding expected utilities $u_a(s,a),u_j(s,a)$, formal and informal care quantities, labor, net public outlay $G(s,a)$, and transition law $P(\cdot\mid s,a)$. These laws are invariant to the government's earlier history and future randomization conditional on $(s,a)$. Every response satisfies household budgets, time and care feasibility, and insurer premium and claim constraints. Need-verification rules have already excluded illegal actions. The government observes $s$ before choosing $a$.

This is a substantial restriction. One sufficient setting has static family bargaining over current care, fixed insurance contracts, and exogenous transitions conditional on current care; discretionary saving and future-policy anticipation are absent. In a genuinely forward-looking model, an augmented state of promised continuation utilities and enforceable contracts may supply such a reduction, but promise keeping and incentive compatibility would have to be established separately. Neither construction is presumed for the general catalogue model.

Let $0<\beta<1$ and bounded rewards
\[
 r(s,a)=w_a u_a(s,a)+\sum_jw_j u_j(s,a)-\lambda G(s,a),
 \qquad \lambda\geq0.
\]
Utilities include consumption and the opportunity cost of caregiver time exactly once. Public outlay $G$ is net of specified receipts. Suppose an outside fiscal account commits an initial present-value amount $B\in\R$ and the legal financing constraint is
\[
 \E\sum_{t\geq0}\beta^tG(s_t,a_t)\leq B.
\]
Thus this is expected present-value financing with access to the stated account; it does not promise solvency in every realized history. The multiplier of this hard constraint below is distinct from the normative cost $\lambda$ already in $r$.

\section{An exact policy characterization}
For any nonanticipative, possibly history-dependent randomized public policy, define its normalized occupation measure
\[
 x(s,a)=(1-\beta)\E\sum_{t\geq0}\beta^t
                         \Ind\{s_t=s,a_t=a\}.
\]

\begin{theorem}[Financed policy equivalence]\label{occupation}
Under the response reduction assumption, the attainable occupation measures are exactly the nonnegative arrays satisfying
\[
 \sum_{a\in\mathcal A(s)}x(s,a)
  =(1-\beta)\mu(s)+\beta\sum_{z,b}x(z,b)P(s\mid z,b)
  \quad(s\in\mathcal S).
\]
Every such array is realized by a stationary randomized public policy. With expected financing, the optimal value is the linear program
\[
 \max_x\ \frac1{1-\beta}\sum_{s,a}r(s,a)x(s,a),\qquad
 \sum_{s,a}G(s,a)x(s,a)\leq(1-\beta)B,
\]
subject to these flow equations. The feasible set is either empty or compact; in the latter case an optimal policy exists.
\end{theorem}
\begin{proof}
For any policy, summing discounted probabilities of visiting state $s$ separates the initial date from later dates. Conditional on state and action, the next-state law is $P$ by assumption, including for history-dependent choices. This gives the flow equations. Absolute convergence follows from bounded terms and discounting.

Conversely, let $x$ satisfy the equations and put $y(s)=\sum_ax(s,a)$. If $y(s)>0$, choose action $a$ with probability $\pi(a\mid s)=x(s,a)/y(s)$; if $y(s)=0$, choose any legal action. The equations become
$y=(1-\beta)\mu+\beta P_\pi^\top y$ in column-vector notation. This linear equation has a unique solution: if $z$ is the difference of two solutions, then
$\|z\|_1=\beta\|P_\pi^\top z\|_1\leq\beta\|z\|_1$, hence $z=0$. The actual normalized discounted occupancy of the stationary policy solves this same equation, so equals $y$. Its state-action occupancy equals $x$ at positive-$y$ states and is zero at all others.

Summing all flow equations gives $\sum_{s,a}x(s,a)=1$. Therefore their nonnegative solution set lies in a finite simplex and is closed. Intersecting with the financing half-space yields a compact set. Discounted utility and expenditure are the stated linear functionals of $x$, proving the exact reduction. The linear objective attains its maximum on a nonempty compact feasible set.
\end{proof}

\begin{theorem}[At most one state needs randomization]\label{one}
If the financing problem is feasible, it has an optimal stationary policy which randomizes between at most two actions in at most one state with positive occupation. In every other occupied state it is deterministic. If the financing inequality is slack at a chosen optimal extreme point, its policy is deterministic at every occupied state.
\end{theorem}
\begin{proof}
A linear objective on a nonempty compact polytope has an optimal extreme point: its optimal face is itself a nonempty compact polytope and has a vertex. Let $x$ be such an optimal vertex, $S_+=\{s:y(s)>0\}$, and $n_+=|S_+|$. For a state outside $S_+$, its flow equation has zero left side; nonnegativity forces $\mu(s)=0$ and $P(s\mid z,b)=0$ for every positive coordinate $x(z,b)$. Hence, restricted to positive coordinates, all flow rows outside $S_+$ vanish.

Let $k$ be the number of positive state-action coordinates. Restricted to these coordinates there are at most $n_+$ flow equalities and at most one binding financing equality. If $k>n_++1$, a nonzero vector in the null space of these rows exists. A sufficiently small positive or negative multiple preserves nonnegativity on the positive coordinates, leaves zero coordinates unchanged, preserves flows and any binding budget, and preserves a slack budget by continuity. This expresses $x$ as the midpoint of two distinct feasible points, contrary to extremality. Therefore $k\leq n_++1$. Since every occupied state needs at least one positive action, only one extra positive action is possible. If the budget is slack, the same null-space argument uses only $n_+$ equalities, giving $k\leq n_+$. The stationary construction in Theorem~\ref{occupation} gives the stated policies.
\end{proof}

This result is useful for implementing a constrained reform: one financing constraint does not require independent lotteries over many cash and in-kind packages in every need state. It does not remove the need to validate each package's family response. Several fiscal or equity constraints can require more randomization.

\section{A checkable upper bound without solving the optimum}
\begin{proposition}[Financing and continuation certificate]
Suppose a vector $v(s)$ and $\eta\geq0$ obey, for every legal $(s,a)$,
\[
 v(s)\geq r(s,a)-\eta G(s,a)
                  +\beta\sum_zP(z\mid s,a)v(z).
\]
Then every financed policy has welfare at most
$\sum_s\mu(s)v(s)+\eta B$. If a feasible occupation measure achieves equality in these inequalities on its support and satisfies
$\eta\{B-(1-\beta)^{-1}\sum Gx\}=0$, it is optimal.
\end{proposition}
\begin{proof}
Multiply each inequality by $x(s,a)$, sum, and use the flow equations to cancel continuation terms. The result is
\[
 (1-\beta)\sum_s\mu(s)v(s)
   \geq\sum_{s,a}r(s,a)x(s,a)-\eta\sum_{s,a}G(s,a)x(s,a).
\]
Divide by $1-\beta$ and use the financing constraint and $\eta\geq0$. If equality holds on occupied actions, summing introduces no slack; the complementary financing condition removes the remaining slack. The feasible policy then attains a bound applying to all policies and is optimal.
\end{proof}

This certificate separates a chosen welfare penalty $\lambda$ from the scarcity value $\eta$ of the hard financing account. It is also valid when net public outlays can be negative in some states. No claim that every state must balance its own budget has been introduced.

\section{Care responses are not production substitution parameters}
Consider one need state, one caregiver, and a current family objective
\[
 U(f,h)=af+bh-\tfrac12(f^2+h^2)+\gamma fh-pf-wh,
 \qquad |\gamma|<1.
\]
Here $f,h$ are care inputs, $p,w$ their effective prices, and positive solutions are considered. The Hessian is negative definite, so the interior optimum is unique. Its first-order conditions are
$f-\gamma h=a-p$ and $h-\gamma f=b-w$.

\begin{proposition}[One reform path does not identify complementarity]
A given pair of before/after positive care choices $(f_0,h_0)$ and $(f_1,h_1)$ can be generated both with $\gamma>0$ and with $\gamma<0$ if the reform's effective price and preference/need shifts are unrestricted. Hence observed joint changes in care do not identify the sign of $U_{fh}$.
If only $p$ changes while $a,b,w,\gamma$ are fixed, then
\[
 \frac{\partial h}{\partial p}=-\frac{\gamma}{1-\gamma^2},
 \qquad \frac{\partial f}{\partial p}=-\frac1{1-\gamma^2},
\]
so the ratio of these responses identifies $\gamma$ in this model.
\end{proposition}
\begin{proof}
For either chosen $\gamma\in(-1,1)$ and each observed pair, choose any nonnegative prices $p_t,w_t$ and set
$a_t=f_t-\gamma h_t+p_t$, $b_t=h_t-\gamma f_t+w_t$.
The observed pair satisfies the first-order conditions, and strict concavity makes it the unique optimum. The construction permits both signs of $\gamma$, proving nonidentification with unrestricted shifts. When only $p$ varies, differentiate the two first-order conditions to get
$f_p-\gamma h_p=-1$ and $h_p-\gamma f_p=0$.
Solve the two equations to obtain the derivatives and their ratio.
\end{proof}

A reform that changes caregiver compensation, insurance coverage, and eligibility simultaneously need not satisfy that exclusion. Nor does the utility cross derivative automatically equal a physical care-production cross derivative. Family bargaining and latent need add further ambiguities.

\section{Unresolved target}
The occupation theorem solves a finite policy problem once a valid response reduction is supplied. General forward-looking saving, private insurance purchase, endogenous premiums, family bargaining, hidden needs, and ambiguity over equilibrium selection remain missing. The observed-data law need not identify $P,r,G$ for untried packages. Welfare bounds over compatible models require keeping each model's joint response and finance restrictions, rather than selecting its most favorable entries independently. The canonical joint family--insurance--public-finance target therefore remains unresolved.

\section{Completion Estimate}
45\%

This is a post hoc, heuristic estimate for the full catalogue target.
A financed finite response model is solved exactly through occupation measures, including a limited randomization theorem and a checkable optimality certificate. Forward looking saving, insurance purchase and premiums, family bargaining, latent needs, and welfare bounds across jointly compatible response models remain outside this reduction.

\begin{thebibliography}{9}
\bibitem{brown} J. R. Brown and A. Finkelstein (2008), \emph{The Interaction of Public and Private Insurance: Medicaid and the Long-Term Care Insurance Market}, American Economic Review 98(3), 1083--1102. Publisher abstract inspected 9 October 2026; full model was not verified here. \url{https://www.aeaweb.org/articles?id=10.1257/aer.98.3.1083}.
\bibitem{mommaerts} C. Mommaerts (2025), \emph{Long-Term Care Insurance and the Family}, Journal of Political Economy 133(1), 1--52. Publisher abstract and publication information inspected 9 October 2026. Online 7 November 2024; January 2025 issue. Full text was access-limited. \url{https://www.journals.uchicago.edu/doi/10.1086/732887}.
\end{thebibliography}

\end{document}
